\documentclass[14pt,a4paper,oneside]{extbook}
\usepackage[utf8]{vietnam}
\usepackage{amsmath,amssymb,mathrsfs,fancyhdr,enumerate,multirow,makecell}
\usepackage{tkz-euclide,tikz-3dplot,tikz,tkz-tab}
\usetikzlibrary{arrows.meta,calc,intersections,angles,quotes,backgrounds,shapes.geometric,patterns,through,shadows,positioning}
\usepackage{yhmath}
\usepackage[top=0.7cm, bottom=0.7cm, left=1.0cm, right=1.0cm, headheight=16pt, headsep=3mm, footskip=5mm]{geometry}
\usepackage[hidelinks,unicode]{hyperref}
\usepackage{qrcode}
\usepackage[many]{tcolorbox}
\usepackage{fontawesome}

\pagestyle{fancy}
\fancyhf{}
\fancyhead[L]{\footnotesize\bfseries\color{teal!85!black}\faBook~SỔ TAY KIẾN THỨC TOÁN 5}
\fancyhead[R]{\footnotesize\bfseries\color{blue!85!black}Lớp Toán Thầy Lĩnh -- loptoanthaylinh.vn}
\fancyfoot[C]{\footnotesize\bfseries\color{gray!80!black}-- Sổ tay kiến thức: Bài 24. Luyện tập chung phép tính số thập phân --}
\renewcommand{\headrulewidth}{0.4pt}

\begin{document}

% --- BANNER TIÊU ĐỀ BÀI HỌC ---
\noindent
\begin{tcolorbox}[
  enhanced,
  colback=teal!8!white,
  colframe=teal!80!black,
  arc=3mm,
  boxrule=1.1pt,
  width=\linewidth,
  halign=center,
  top=2.2mm, bottom=2.2mm,
  left=3mm, right=3mm,
  drop fuzzy shadow
]
  {\large\bfseries\color{teal!95!black}\faBookmark\ BÀI 24. LUYỆN TẬP CHUNG PHÉP TÍNH SỐ THẬP PHÂN}\\[2pt]
  {\scriptsize\color{black!75}\itshape Thứ tự biểu thức, kỹ thuật đặt thừa số chung và giải toán đố số thập phân trên 1 trang A4}
\end{tcolorbox}

\vspace{3.2mm}
% --- KHỐI 1: BẢNG QUY TẮC BỐN PHÉP TÍNH STP ---
\noindent
\begin{tcolorbox}[
  enhanced,
  title={\footnotesize\bfseries\color{white}\faCalculator\ 1. Tổng kết nhanh quy tắc 4 phép tính với số thập phân},
  colback=blue!2!white,
  colframe=blue!75!black,
  arc=2.5mm, boxrule=0.8pt, top=1.8mm, bottom=1.8mm, left=2.5mm, right=2.5mm
]
  \scriptsize
  \centering
  \begin{tabular}{|l|l|l|}
    \hline
    \textbf{Phép tính} & \textbf{Quy tắc đặt tính} & \textbf{Xử lý dấu phẩy} \\ \hline
    \textbf{Cộng / Trừ} & Thẳng cột các hàng và dấu phẩy & Hạ dấu phẩy ở kết quả thẳng cột với các dấu phẩy trên \\ \hline
    \textbf{Nhân} & Nhân như số tự nhiên (không cần thẳng phẩy) & Đếm tổng chữ số TP của 2 thừa số $\rightarrow$ Tách ở tích từ phải sang trái \\ \hline
    \textbf{Chia} & Bỏ dấu phẩy ở SC, dịch phẩy ở SBC tương ứng & Đặt dấu phẩy ở thương ngay khi dùng đến chữ số đầu tiên của phần TP \\ \hline
  \end{tabular}
\end{tcolorbox}

\vspace{3mm}

% --- KHỐI 2: KỸ THUẬT TÍNH THUẬN TIỆN ĐỈNH CAO ---
\noindent
\begin{minipage}[t]{0.52\textwidth}
  \begin{tcolorbox}[
    enhanced,
    title={\footnotesize\bfseries\color{white}\faBolt\ 2. Kỹ thuật đưa về cùng thừa số chung},
    colback=teal!2!white,
    colframe=teal!75!black,
    arc=2.5mm, boxrule=0.8pt, top=1.8mm, bottom=1.8mm, left=2.2mm, right=2.2mm
  ]
    \footnotesize
    \textit{VD 1:} $4,8 \times 2,5 + 4,8 \times 7,5$\\
    $= 4,8 \times (2,5 + 7,5) = 4,8 \times 10 = \mathbf{48}$.\\[2pt]
    \textit{VD 2:} $9,64 \times 101 - 9,64$\\
    $= 9,64 \times (101 - 1) = 9,64 \times 100 = \mathbf{964}$.\\[2pt]
    \textit{VD 3 (Đổi thừa số):} $1,25 \times 0,7 \times 8$\\
    $= (1,25 \times 8) \times 0,7 = 10 \times 0,7 = \mathbf{7}$.
  \end{tcolorbox}
\end{minipage}%
\hfill
\begin{minipage}[t]{0.46\textwidth}
  \begin{tcolorbox}[
    enhanced,
    title={\footnotesize\bfseries\color{white}\faBalanceScale\ 3. Thứ tự tính biểu thức},
    colback=purple!2!white,
    colframe=purple!70!black,
    arc=2.5mm, boxrule=0.8pt, top=1.8mm, bottom=1.8mm, left=2mm, right=2mm
  ]
    \footnotesize
    $\bullet$ Trong ngoặc trước, ngoài ngoặc sau: $( \ ) \rightarrow [ \ ]$.\\
    $\bullet$ Nhân, chia trước $\rightarrow$ Cộng, trừ sau.\\
    $\bullet$ Chỉ gồm phép nhân, chia (hoặc chỉ có cộng, trừ): Tính từ \textbf{trái sang phải}.
  \end{tcolorbox}
\end{minipage}

\vspace{3mm}

% --- KHỐI 3: SƠ ĐỒ TƯ DUY ---
\noindent
\begin{tcolorbox}[
  enhanced,
  title={\footnotesize\bfseries\color{white}\faSitemap\ Sơ đồ tư duy: Bốn phép tính số thập phân},
  colback=yellow!2!white,
  colframe=teal!85!black,
  arc=3mm, boxrule=0.9pt, top=1.8mm, bottom=1.8mm, left=1.5mm, right=1.5mm,
  halign=center
]
  \begin{tikzpicture}[
    scale=0.92,
    every node/.style={transform shape},
    root/.style={rectangle, rounded corners=5pt, fill=teal!85!black, text=white, font=\bfseries\footnotesize, inner sep=5pt},
    b1/.style={rectangle, rounded corners=4pt, fill=blue!12, draw=blue!75!black, text=blue!90!black, font=\bfseries\scriptsize, inner sep=4pt},
    b2/.style={rectangle, rounded corners=4pt, fill=red!12, draw=red!75!black, text=red!90!black, font=\bfseries\scriptsize, inner sep=4pt},
    b3/.style={rectangle, rounded corners=4pt, fill=orange!12, draw=orange!75!black, text=orange!90!black, font=\bfseries\scriptsize, inner sep=4pt},
    b4/.style={rectangle, rounded corners=4pt, fill=green!15, draw=green!65!black, text=green!80!black, font=\bfseries\scriptsize, inner sep=4pt}
  ]
    \node[root] (core) at (0,0) {\faCalculator\ TÍNH TOÁN STP};
    \node[b1] (n1) at (-4.7, 0.8) {CỘNG: Thẳng cột phẩy};
    \node[b2] (n2) at (-4.7, -0.8) {TRỪ: Thêm 0 nếu thiếu};
    \node[b3] (n3) at (4.7, 0.8) {NHÂN: Đếm tổng chữ số TP};
    \node[b4] (n4) at (4.7, -0.8) {CHIA: Bỏ phẩy ở số chia};

    \draw[teal!85!black, thick, -{Stealth[length=2mm]}] (core) -- (n1);
    \draw[teal!85!black, thick, -{Stealth[length=2mm]}] (core) -- (n2);
    \draw[teal!85!black, thick, -{Stealth[length=2mm]}] (core) -- (n3);
    \draw[teal!85!black, thick, -{Stealth[length=2mm]}] (core) -- (n4);
  \end{tikzpicture}
\end{tcolorbox}

\vspace{3mm}

% --- KHỐI 4: TÌM X PHỨC HỢP & BẪY TRẮC NGHIỆM ---
\noindent
\begin{minipage}[t]{0.54\textwidth}
  \begin{tcolorbox}[
    enhanced,
    title={\footnotesize\bfseries\color{white}\faSearch\ Dạng toán tìm $x$ chứa phép tính gộp},
    colback=green!3!white,
    colframe=green!60!black,
    arc=2.5mm, boxrule=0.8pt, top=1.8mm, bottom=1.8mm, left=2.2mm, right=2.2mm
  ]
    \footnotesize
    $x \times 2,5 + x \times 7,5 = 35$\\
    $\triangleright$ Đặt thừa số chung: $x \times (2,5 + 7,5) = 35$\\
    $\triangleright$ $x \times 10 = 35 \Rightarrow x = 35 : 10 \Rightarrow \mathbf{x = 3,5}$.\\
    Luôn tìm cách thu gọn biểu thức chứa $x$ trước khi tìm.
  \end{tcolorbox}
\end{minipage}%
\hfill
\begin{minipage}[t]{0.44\textwidth}
  \begin{tcolorbox}[
    enhanced,
    title={\footnotesize\bfseries\color{white}\faExclamationTriangle\ Mẹo tránh bẫy trắc nghiệm},
    colback=red!2!white,
    colframe=red!70!black,
    arc=2.5mm, boxrule=0.8pt, top=1.8mm, bottom=1.8mm, left=2.2mm, right=2.2mm
  ]
    \footnotesize
    $\bullet$ Biểu thức $a \times b - a$: Nhớ rằng số $a$ đứng một mình chính là $a \times \mathbf{1}$.\\
    \textit{VD:} $5,6 \times 99 + 5,6 = 5,6 \times (99 + 1) = 560$.\\
    $\bullet$ Cẩn thận khi nhân chia phân phối với dấu trừ.
  \end{tcolorbox}
\end{minipage}

\end{document}
